% Lösungen Einheit 2 von 4 – Satz vom Nullprodukt (zusammenbau v0.1)
\einheitenkopf{Einheit 2 von 4 · Satz vom Nullprodukt}

\erg{14}{a) gilt – rechts steht null \quad b) $x - 2 = 0$ oder $x - 5 = 0$; $x_1 = 2$, $x_2 = 5$ \quad c) $x + 4 = 0$ oder $x - 1 = 0$; $x_1 = -4$, $x_2 = 1$ \quad d) $x = 0$ oder $x + 6 = 0$; $x_1 = 0$, $x_2 = -6$ \quad e) beide Faktoren gleich, genau eine Lösung: $x = 7$ \quad f) $(-5 + 5) \cdot (-5 - 1) = 0 \cdot (-6) = 0$ (wA): ja \quad g) $x \cdot (x + 9) = 0$; $x_1 = 0$, $x_2 = -9$ \quad h) $x \cdot (x - 3) = 0$; $x_1 = 0$, $x_2 = 3$ \quad i) $x \cdot (2x - 6) = 0$; $x_1 = 0$, $x_2 = 3$ \quad j) $x^2 + 3x = 10$, also $x^2 + 3x - 10 = 0$ \quad k) $x = -7$, denn $(-7) \cdot (-7 + 9) = (-7) \cdot 2 = -14$}
\erg{15}{a) Durch $x$ geteilt, dabei geht $x = 0$ verloren. Richtig: $x^2 - 6x = 0$, $x \cdot (x - 6) = 0$, $x_1 = 0$, $x_2 = 6$ \quad b) Null mal jede Zahl ergibt null – egal, welchen Wert der andere Faktor hat}
